%!TEX program = lualatex
\documentclass[ngerman,headheight=70pt]{scrartcl}
\usepackage[ngerman]{babel}
\usepackage{blindtext}
\usepackage[utf8]{luainputenc}
\usepackage{scrlayer-scrpage}
\usepackage[T1]{fontenc}
\usepackage{fontspec}
\usepackage{multicol}
\usepackage{enumitem}
\usepackage{epigraph}
\usepackage[german=quotes]{csquotes}
\usepackage[bottom=2cm,footskip=8mm]{geometry}
\setmainfont{Source Sans Pro}
\MakeOuterQuote{"}

\setlength{\parindent}{0pt}
\setlength{\parskip}{6pt}

\pagestyle{plain.scrheadings}
\thispagestyle{scrheadings}
\lohead{\vspace*{-1.5cm}\includegraphics[scale=0.4]{../images/stupa_logo_rot.pdf}}
\rohead{\\Präsidium des\\Studierendenparlaments}

\lofoot[Universität Hamburg · Präsidium des Studierendenparlaments\\\hrule Von-Melle-Park 5 · D-20146 Hamburg]{Universität Hamburg · Präsidium des Studierendenparlaments\\\hrule Von-Melle-Park 5 · D-20146 Hamburg}
\cofoot[]{}
\cohead[\thepage]{}

\renewcommand{\thesubsection}{\arabic{section}.\arabic{subsection}}
\renewcommand{\thesubsubsection}{\arabic{subsubsection}.}
\setcounter{subsection}{-1}

\begin{document}
    UHH · StuPa-Präsidium · Von-Melle-Park 5 · D-20146 Hamburg

    \section{Inhaltliches Konzept für das Wintersemester}

    \subsection{StuPa-Sitzungen}

    Für die erste StuPa-Sitzung am 27. Oktober schlagen wir als Präsidium die
    Themen "Wahlbeteiligung bei der StuPa-Wahl" und "Haushaltsberatungen in
    der Bürgerschaft" als TOPs vor.

    Hintergrund ist die geringe Wahlbeteiligung von rund 13\% bei der letzten
    Wahl. Wir möchten Gelegenheit schaffen, Ideen zu sammeln, was zur Erhöhung
    der Wahlbeteiligung gemacht werden kann. Jede*r ist herzlich eingeladen, sich
    daran zu beteiligen.

    Das zweite Thema ergibt sich aus dem Fakt, dass zu dem Zeitpunkt der nächste
    Haushalt der Stadt in den Fachausschüssen beraten wird und Anfang November
    der Haushaltsausschuss die Stellungnahmen zusammentragen wird. Die
    Oktobersitzung des Studierendenparlaments ist daher die einzige Möglichkeit,
    um vorher noch einen Antrag zu verabschieden und den Abgeordneten zukommen
    zu lassen, der eine bedarfsgerechte Ausfinanzierung der Uni fordert.
    Angesichts der Tatsache, dass eine Unterfinanzierung die Universität als
    Ganzes bedroht und bereits einige Studiengänge gefährdet sind oder sogar
    bereits nicht mehr angeboten werden, wäre es daher wünschenswert, wenn die
    Verfasste Studierendenschaft gemeinsam für eine Verbesserung der Situation
    eintritt.

    Darüber hinaus schlagen wir vor, am 15. Dezember im Rahmen einer Aktuellen
    Stunde über die Haushaltsberatungen zu sprechen, denn am 13., 14. und 15.
    Dezember sind die  Haushaltsberatungen in der Bürgerschaft, an deren Ende der
    Haushalt verabschiedet wird.


    \subsection{Wahlmotto}

    Nach reiflicher Überlegung sind wir zu folgendem Wahlmotto für die diesjährige
    Wahl des Studierendenparlaments gekommen, welches wir dem Parlament vorschlagen
    möchten:

    "Fremdbestimmt oder selbstgestaltet? Universität gemeinsam entwickeln!"

    Für Kommentare und Kritik, weitere Ideen oder Ergänzungen sind wir offen.

    \subsection{Wahlveranstaltungen}
    Noch vor Beginn der heißen Wahlkampfphase wollen wir eine inhaltliche
    Veranstaltung zum Thema Studienreform veranstalten. Insbesondere soll dort
    über die Frage "Wie soll Studium aussehen?" diskutiert werden. Damit wollen
    wir Studierende ermuntern, ihre Vorstellungen von einer guten Forschung,
    Lehre und Studium zu artikulieren und zugleich Interesse wecken und uns
    gegenseitig qualifizieren für den Dies Academicus am 30. November. Die
    Veranstaltung wird voraussichtlich am 10. November vor der StuPa-Sitzung
    stattfinden.

    Vor der StuPa-Sitzung am 15. Dezember wird es die erste Listenvorstellung
    geben, wobei das Konzept dieses Jahr etwas anders ist. Statt "Werbeblöcke"
    soll es stärker um die Herausarbeitung auch der weltanschaulichen Positionen
    der Listen gehen. Dabei geht es um grundsätzliche Fragen: "Wie soll Hamburg
    in zehn Jahren aussehen?", "Wie soll die Gesellschaft aussehen?" und "Wie
    soll Hochschule aussehen?" Das Ziel dieser ersten Vorstellung ist es, über
    die grundsätzlichen Ansichten der Hochschulgruppen zu reden und dies mit
    den aktuellen Fragen der Zeit zu verbinden.
    Wie in den vergangenen Jahren soll es die Möglichkeit geben, die Liste im
    Vorfeld der Veranstaltung in Form von Infoständen vorzustellen.

    Vor der StuPa-Sitzung am 12. Januar gibt es die zweite Listenvorstellung,
    bei der es um die Frage "Was wollt ihr in einem Jahr in der Verfassten
    Studierendenschaft machen?" gehen wird. Hier kann es dann konkreter um
     hochschulpolitische Aktivitäten und AStA-Arbeit gehen.
    Wir gehen davon aus, dass hier schon so viele eigenständige
    Wahlkampfaktivitäten der Gruppen stattfinden, dass wir keinen "Markt der
    Möglichkeiten" brauchen.

    Insgesamt hoffen wir, mit den beiden Vorstellungen potentiellen Wähler*innen
    wie den Hochschulaktiven die Möglichkeit zu eröffnen, aus dem Studien- oder
    Gremienalltag aus- und in eine viel zu selten geführte Diskussion über
    Visionen und Motive des politischen Engagements einzutreten. Damit möge auch
    die Bedeutung der Wahlen zum Studierendenparlament (und Akademischen Senat)
    deutlicher werden.
\end{document}
